% Source: physical PDF pp. 310--328 (printed pp. 287--305), from the
% Section 29.5 heading through the final paragraph immediately before
% Problems.

\subsection[The Seiberg--Witten Solution]{The Seiberg--Witten
Solution\footnotemark}
\footnotetext{This section lies somewhat out of the book's main line of
development and may be omitted in a first reading.}

It often happens that the tree-approximation potential in a supersymmetric
theory will take the value zero for a continuous range of scalar field
values. (For one example, see Eq.~\eqref{eq:29.1.9}.) In this case the theory
has a number of scalar excitations with zero mass in the tree approximation,
which since supersymmetry is unbroken must be accompanied with suitable
fermionic superpartners. At low energies the theory will then be described by
a family of supersymmetric effective Lagrangians, whose members are
parameterized by one or more \textit{moduli}, the scalar expectation values
in the underlying theory. Quantum effects in the underlying theory can
modify the dependence of the effective Lagrangian on these moduli, and even
alter the topology of the space of moduli.~\hyperref[ch29-ref-9]{[9]}

In one of the most striking accomplishments of the 1990s in supersymmetry
theory, Seiberg and Witten~\hyperref[ch29-ref-10]{[10]} were able to
calculate the exact dependence of the low-energy effective Lagrangian on a
modular parameter in gauge theories with \(N=2\) supersymmetry. The ideas
behind this calculation can be made apparent by running through only the
simplest special case, that of an \(SU(2)\) gauge theory with \(N=2\)
supersymmetry and no additional matter hypermultiplets.

We saw in Section 27.9 that the Lagrangian density for this theory is given
after elimination of the auxiliary fields by
\begin{equation}
  \begin{aligned}
  \mathcal{L}
  =\frac{1}{e^2}\Biggl[
  &-\sum_A(D_\mu\phi)_A^*(D^\mu\phi)_A
  -\frac12\sum_A\left(\bar\psi_A(\sl{D}\psi)_A\right)
  \\
  &-2\sqrt{2}\operatorname{Re}\sum_{ABC}
    \epsilon_{ABC}
    \left(\lambda_{AL}^{\mathrm T}\epsilon\psi_{CL}\right)\phi_B^*
  -\frac14\sum_A f_{A\mu\nu}f_A^{\mu\nu}
  \\
  &-\frac12\sum_A\left(\bar\lambda_A(\sl{D}\lambda)_A\right)
  \Biggr]
  +\frac{\theta}{64\pi^2}\epsilon_{\mu\nu\rho\sigma}
    \sum_A f_A^{\mu\nu}f_A^{\rho\sigma}
  -V(\phi,\phi^*)\,.
  \end{aligned}
  \tag{29.5.1}\label{eq:29.5.1}
\end{equation}
Here \(A\), \(B\), and \(C\) now run over the values 1, 2, and 3. We have now
rescaled all the fields by multiplying each of them by a factor \(e\), so
that \(e\) does not appear in the covariant derivatives:
\begin{equation}
  (D_\mu\psi)_A
  =\partial_\mu\psi_A+\sum_{BC}\epsilon_{ABC}V_{B\mu}\psi_C\,,
  \tag{29.5.2}\label{eq:29.5.2}
\end{equation}
\begin{equation}
  (D_\mu\lambda)_A
  =\partial_\mu\lambda_A+\sum_{BC}\epsilon_{ABC}V_{B\mu}\lambda_C\,,
  \tag{29.5.3}\label{eq:29.5.3}
\end{equation}
\begin{equation}
  (D_\mu\phi)_A
  =\partial_\mu\phi_A+\sum_{BC}\epsilon_{ABC}V_{B\mu}\phi_C\,,
  \tag{29.5.4}\label{eq:29.5.4}
\end{equation}
\begin{equation}
  f_{A\mu\nu}
  =\partial_\mu V_{A\nu}-\partial_\nu V_{A\mu}
  +\sum_{BC}\epsilon_{ABC}V_{B\mu}V_{C\nu}\,,
  \tag{29.5.5}\label{eq:29.5.5}
\end{equation}
and the potential is
\begin{equation}
  V(\phi,\phi^*)
  =2\sum_A\left[
    \sum_{BC}\epsilon_{ABC}\operatorname{Re}\phi_B
    \operatorname{Im}\phi_C
  \right]^2.
  \tag{29.5.6}\label{eq:29.5.6}
\end{equation}
This potential takes the value zero for a family of scalar field expectation
values, which (up to a gauge transformation) may be parameterized as
\begin{equation}
  \phi_1=\phi_2=0\,,
  \qquad
  \phi_3=a\,,
  \tag{29.5.7}\label{eq:29.5.7}
\end{equation}
with \(a\) a complex parameter, known as the vacuum modulus. This vacuum
expectation value gives masses \(2\lvert a\rvert\) to the vector fields
\(V_{1\mu}\) and \(V_{2\mu}\), the gauginos \(\lambda_1\) and \(\lambda_2\),
the chiral fermions \(\psi_1\) and \(\psi_2\), and the scalars \(\phi_1\) and
\(\phi_2\), leaving \(V_{3\mu}\), \(\lambda_3\), \(\psi_3\), and
\(\phi_3=a\) all massless.

Taking account only of these zero-mass modes (and dropping the subscript 3),
the tree approximation gives the effective low-energy theory as simply the
free-field theory with Lagrangian density
\begin{equation}
  \begin{aligned}
  \mathcal{L}_{\mathrm{eff}}
  =\frac{1}{e^2}\Biggl[
  &-(\partial_\mu a)^*(\partial^\mu a)
  -\frac12\left(\bar\psi(\sl{\partial}\psi)\right)
  -\frac14 f_{\mu\nu}f^{\mu\nu}
  \\
  &-\frac12\left(\bar\lambda(\sl{\partial}\lambda)\right)
  \Biggr]
  +\frac{\theta}{64\pi^2}
    \epsilon_{\mu\nu\rho\sigma}f^{\mu\nu}f^{\rho\sigma}\,,
  \end{aligned}
  \tag{29.5.8}\label{eq:29.5.8}
\end{equation}
with \(f_{\mu\nu}=\partial_\mu V_\nu-\partial_\nu V_\mu\). Indeed, any
renormalizable theory of a single gauge boson and its gauge-neutral \(N=2\)
superpartners must be a free-field theory, because \(N=2\) supersymmetry
does not allow a superpotential for this theory.

But this is not the whole story. In integrating out the massive degrees of
freedom in this theory, quantum corrections produce non-renormalizable
interaction terms in the low-energy effective field theory. We can classify
the interactions and Feynman diagrams containing them that predominate at
low energy by the same sort of counting of powers of energy that we used in
dealing with low-momentum pions and nucleons in Section 19.5. When we use the
effective Lagrangian perturbatively to calculate low-energy scattering
amplitudes, the number \(\nu\) of powers of energy contributed by a connected
graph with \(L\) loops, \(I_f\) external fermion lines, \(I_b\) internal boson
\((a\) or \(V_\mu)\) lines, \(I_a\) internal auxiliary field lines, and
\(V_i\) vertices of each type \(i\), is
\begin{equation}
  \nu=4L+\sum_iV_id_i-2I_b-I_f\,,
  \tag{29.5.9}\label{eq:29.5.9}
\end{equation}
where \(d_i\) is the number of derivatives in the interaction of type \(i\).
(Internal auxiliary field lines do not contribute in Eq.~\eqref{eq:29.5.9}
because their propagators are independent of momentum.) These quantities
are subject to the topological relations
\begin{equation}
  L=I_b+I_f+I_a-\sum_iV_i+1\,,
  \tag{29.5.10}\label{eq:29.5.10}
\end{equation}
and
\begin{equation}
  2I_b+E_b=\sum_iV_ib_i\,,
  \qquad
  2I_f+E_f=\sum_iV_if_i\,,
  \qquad
  2I_a+E_a=\sum_iV_ia_i\,,
  \tag{29.5.11}\label{eq:29.5.11}
\end{equation}
where \(E_b\), \(E_f\), and \(E_a\) are the numbers of external boson,
fermion, and auxiliary field lines, and \(b_i\), \(f_i\), and \(a_i\) are the
numbers of boson, fermion, and auxiliary fields in interactions of type
\(i\). We can therefore write the number of powers of energy as
\begin{equation}
  \nu=\sum_iV_i\left(d_i+\frac12f_i+a_i-2\right)
  +2L-E_f-2E_a+2\,.
  \tag{29.5.12}\label{eq:29.5.12}
\end{equation}
According to Eqs.~\eqref{eq:26.8.4} and \eqref{eq:27.4.42}, both the
\(D\)-term of a function of left-chiral scalar \(N=1\) superfields and
their adjoints and the \(\mathcal{F}\)-term of a pair of \(N=1\) gauge
superfields \(W_\alpha\) times an arbitrary function of \(N=1\) left-chiral
scalar superfields have \(d_i+\tfrac12f_i+a_i=2\), while adding any
additional \(W_\alpha\) factors or superderivatives \(\mathcal{D}_\alpha\)
would give \(d_i+\tfrac12f_i+a_i>2\), so in our case supersymmetry rules out
any interactions with \(d_i+\tfrac12f_i+a_i\) less than 2. The dominant
contribution to low-energy scattering amplitudes is thus given by the tree
approximation (\(L=0\)), calculated with an effective Lagrangian containing
only terms with \(d_i+\tfrac12f_i+a_i=2\), which takes the general form
discussed in Section 27.4:
\begin{equation}
  \mathcal{L}_{\mathrm{eff}}
  =\frac12\left[K(\Phi,\Phi^*)\right]_D
  -\frac12\operatorname{Re}
    \left[T(\Phi)\left(W_L^{\mathrm T}\epsilon W_L\right)\right]_{\mathcal F}.
  \tag{29.5.13}\label{eq:29.5.13}
\end{equation}
Eq.~\eqref{eq:27.4.42} then gives the Lagrangian in terms of component
fields as
\begin{equation}
  \begin{aligned}
  \mathcal{L}
  ={}&
  \frac{\partial^2K(a,a^*)}{\partial a\,\partial a^*}
  \left[
    -\frac12\left(\bar\psi\sl{\partial}\psi\right)
    +\lvert\mathcal{F}\rvert^2
    -\partial_\mu a\,\partial^\mu a^*
  \right]
  \\
  &-\operatorname{Re}\left\{
    \frac{\partial^3K(a,a^*)}{\partial^2a\,\partial a^*}
    \left(\bar\psi\psi_L\right)\mathcal{F}^*
  \right\}
  -\frac12\operatorname{Re}\left\{
    \frac{\partial^3K(a,a^*)}{\partial^2a\,\partial a^*}
    \left(\bar\psi\gamma^\mu\gamma_5\psi\right)\partial_\mu a
  \right\}
  \\
  &+\frac14
    \frac{\partial^4K(a,a^*)}{\partial^2a\,\partial^2a^*}
    \left(\bar\psi\psi_L\right)\left(\bar\psi\psi_R\right)
  \\
  &+\frac14\operatorname{Re}\left\{
    \left(\bar\lambda\lambda_L\right)
    \left(\bar\psi\psi_L\right)\frac{d^2T(a)}{da^2}
  \right\}
  -\frac12\operatorname{Re}\left\{
    \left(\bar\lambda\lambda_L\right)\mathcal{F}\frac{dT(a)}{da}
  \right\}
  \\
  &+\operatorname{Re}\left\{
    T(a)\left[
      -\frac12\left(\bar\lambda\sl{\partial}(1-\gamma_5)\lambda\right)
      -\frac14f_{\mu\nu}f^{\mu\nu}
      +\frac{i}{8}\epsilon_{\mu\nu\rho\sigma}
        f^{\mu\nu}f^{\rho\sigma}
      +\frac12D^2
    \right]
  \right\}
  \\
  &+\frac{\sqrt{2}}{4}\operatorname{Re}\left\{
    \frac{dT(a)}{da}\left[
      -\left(\bar\lambda\gamma^\mu\gamma^\nu\psi_L\right)f_{\mu\nu}
      +2i\left(\bar\lambda\psi_L\right)D
    \right]
  \right\}.
  \end{aligned}
  \tag{29.5.14}\label{eq:29.5.14}
\end{equation}
In order to implement \(N=2\) supersymmetry, we now want to impose invariance
under the discrete \(R\)-symmetry transformation \eqref{eq:27.9.2}:
\begin{equation}
  \psi\longrightarrow\lambda\,,
  \qquad
  \lambda\longrightarrow-\psi\,,
  \tag{29.5.15}\label{eq:29.5.15}
\end{equation}
with \(a\) and \(V_\mu\) unchanged. The condition that the coefficients of
\((\bar\psi\sl{\partial}\psi)\) and
\((\bar\lambda\sl{\partial}\lambda)\) should be equal is
\begin{equation}
  \frac{\partial^2K(a,a^*)}{\partial a\,\partial a^*}
  =\operatorname{Re}T(a).
  \tag{29.5.16}\label{eq:29.5.16}
\end{equation}
The right-hand side is the sum of a function of \(a\) and a function of
\(a^*\), so
\(\partial^4K/\partial^2a\,\partial^2a^*=0\), and therefore the term quartic
in \(\psi\), which would have had no counterpart for \(\lambda\), is absent.
By an integration by parts, the term
\(\tfrac12\operatorname{Re}\{T(a)(\bar\lambda\sl{\partial}\gamma_5\lambda)\}\)
may be replaced with
\(-\tfrac14\operatorname{Re}\{(\bar\lambda\gamma^\mu\gamma_5\lambda)
\partial_\mu T(a)\}\). The condition that the coefficients of
\((\bar\psi\gamma^\mu\gamma_5\psi)\) and
\((\bar\lambda\gamma^\mu\gamma_5\lambda)\) should be equal is then satisfied
if
\[
  \frac14\partial_\mu T
  =\frac12
    \frac{\partial^3K}{\partial^2a\,\partial a^*}\partial_\mu a\,,
\]
which is also an automatic consequence of Eq.~\eqref{eq:29.5.16}.
According to Eq.~\eqref{eq:26.A.7}, the terms proportional to
\(f_{\mu\nu}(\bar\lambda[\gamma^\mu,\gamma^\nu]\psi)\) and
\(f_{\mu\nu}(\bar\lambda[\gamma^\mu,\gamma^\nu]\gamma_5\psi)\) are
automatically invariant under the transformation \eqref{eq:29.5.15}, as are
also the terms proportional to
\((\bar\lambda\lambda_L)(\bar\psi\psi_L)\) and its adjoint. On the other
hand, the invariance of the terms proportional to
\((\bar\lambda\psi)D\) and \((\bar\lambda\gamma_5\psi)D\) requires that we
extend the transformation \eqref{eq:29.5.15} so that
\begin{equation}
  D\longrightarrow-D\,,
  \tag{29.5.17}\label{eq:29.5.17}
\end{equation}
which also leaves the term \(\tfrac12D^2\) invariant. Finally,
Eq.~\eqref{eq:29.5.16} tells us that the coefficients of
\((\bar\psi\psi_L)\mathcal{F}^*\) and
\((\bar\lambda\lambda_L)\mathcal{F}\) are equal, so the transformation
\eqref{eq:29.5.15} must also be extended so that
\begin{equation}
  \mathcal{F}\longrightarrow\mathcal{F}^*\,,
  \tag{29.5.18}\label{eq:29.5.18}
\end{equation}
which also leaves the term \(\lvert\mathcal{F}\rvert^2\) invariant. We
conclude then that the condition \eqref{eq:29.5.16} ensures that the whole
action derived from Eq.~\eqref{eq:29.5.14} is invariant under the combined
transformation \eqref{eq:29.5.15}, \eqref{eq:29.5.17}, and
\eqref{eq:29.5.18}. Since the Lagrangian is invariant under \(N=1\)
supersymmetry with a left-chiral scalar supermultiplet
\((a,\psi,\mathcal{F})\) and a gauge supermultiplet
\((V_\mu,\lambda,D)\), it is also invariant under a second supersymmetry
with a left-chiral scalar supermultiplet \((a,\lambda,\mathcal{F}^*)\) and a
gauge supermultiplet \((V_\mu,-\psi,-D)\). Eq.~\eqref{eq:29.5.16} is
therefore enough to ensure that the action obtained from
Eq.~\eqref{eq:29.5.13} or \eqref{eq:29.5.14} is invariant (without imposing
the field equations) under \(N=2\) supersymmetry.

The general solution of Eq.~\eqref{eq:29.5.16} can be expressed as
\footnote{The factor \(1/4\pi i\) is included in order to simplify the
duality transformation introduced below. Seiberg and Witten introduced a
function \(\mathcal{F}(a)\) (not related to the auxiliary field
\(\mathcal{F}\)), known as the prepotential, for which
\(h(a)=d\mathcal{F}(a)/da\).}
\begin{equation}
  T(a)=\frac{1}{4\pi i}\frac{dh(a)}{da}\,,
  \qquad
  K(a,a^*)=\operatorname{Im}\left\{\frac{a^*h(a)}{4\pi}\right\},
  \tag{29.5.19}\label{eq:29.5.19}
\end{equation}
with \(h\) a function of \(a\) alone. In terms of \(h\), the Lagrangian
density \eqref{eq:29.5.14} now reads
\begin{equation}
  \begin{aligned}
  \mathcal{L}
  ={}&\frac{1}{4\pi}\operatorname{Im}\Biggl\{
    \left[\frac{dh}{da}\right]\Biggl[
      -\frac12\left(\bar\psi\sl{\partial}(1-\gamma_5)\psi\right)
      -\frac12\left(\bar\lambda\sl{\partial}(1-\gamma_5)\lambda\right)
  \\
  &\hspace{31mm}
      -\partial_\mu a\,\partial^\mu a^*
      +\lvert\mathcal{F}\rvert^2+\frac12D^2
      -\frac14f_{\mu\nu}f^{\mu\nu}
      +\frac{i}{8}\epsilon_{\mu\nu\rho\sigma}
        f^{\mu\nu}f^{\rho\sigma}
    \Biggr]\Biggr\}
  \\
  &+\frac{1}{4\pi}\operatorname{Im}\Biggl\{
    \left[\frac{d^2h}{da^2}\right]\Biggl[
      -\frac12\left(\bar\psi\psi_L\right)\mathcal{F}^*
      -\frac12\left(\bar\lambda\lambda_L\right)\mathcal{F}
      +\frac{\sqrt{2}}{2}i\left(\bar\lambda\psi_L\right)D
  \\
  &\hspace{43mm}
      -\frac{\sqrt{2}}{4}
        \left(\bar\lambda\gamma^\mu\gamma^\nu\psi_L\right)f_{\mu\nu}
    \Biggr]\Biggr\}
  \\
  &+\frac{1}{16\pi}\operatorname{Im}\left\{
    \frac{d^3h(a)}{da^3}
    \left(\bar\lambda\lambda_L\right)\left(\bar\psi\psi_L\right)
  \right\}.
  \end{aligned}
  \tag{29.5.20}\label{eq:29.5.20}
\end{equation}
This is invariant under an \(SU(2)\) \(R\)-symmetry, under which
\((\psi,\lambda)\) transforms as a doublet and
\((\operatorname{Im}\mathcal{F},\operatorname{Re}\mathcal{F},D/\sqrt{2})\)
as a triplet. The discrete transformation \eqref{eq:29.5.15},
\eqref{eq:29.5.17}, and \eqref{eq:29.5.18} is a finite element of this
\(SU(2)\) group: a rotation of \(\pi\) radians around the 2-axis.

Comparing Eqs.~\eqref{eq:29.5.8} and \eqref{eq:29.5.20}, we see that in the
tree approximation
\begin{equation}
  h(a)_{\mathrm{tree}}
  =\left[\frac{4\pi i}{e^2}+\frac{\theta}{2\pi}\right]a\,.
  \tag{29.5.21}\label{eq:29.5.21}
\end{equation}
The achievement of Seiberg and Witten was to calculate \(h(a)\) exactly.

The first step in this calculation is to recognize that there are various
linear transformations on \(a\) and \(dh(a)/da\) that give physically
equivalent theories. This is because of a remarkable property of the
low-energy effective theory, related to the duality property discussed
briefly in Section 27.9. To demonstrate this property, let us return to the
Lagrangian density \eqref{eq:29.5.13} expressed in terms of \(N=1\)
superfields, now using the relations \eqref{eq:29.5.19} required by \(N=2\)
supersymmetry
\begin{equation}
  \mathcal{L}_{\mathrm{eff}}
  =\frac{1}{8\pi}\operatorname{Im}\left[\Phi^*h(\Phi)\right]_D
  -\frac{1}{8\pi}\operatorname{Im}
    \left[h'(\Phi)\left(W_L^{\mathrm T}\epsilon W_L\right)\right]_{\mathcal F}.
  \tag{29.5.22}\label{eq:29.5.22}
\end{equation}
In path integrals the spinor field-strength superfield \(W_L\) is constrained
by the supersymmetric extension \eqref{eq:27.2.20} of the homogeneous Maxwell
equations:
\begin{equation}
  \operatorname{Re}\left(\mathcal{D}_L^{\mathrm T}\epsilon W_L\right)=0.
  \tag{29.5.23}\label{eq:29.5.23}
\end{equation}
This is usually imposed by requiring that \(W_L\) take the form
\eqref{eq:27.2.15}
\begin{equation}
  W_L=\frac{i}{4}
  \left(\mathcal{D}_R^{\mathrm T}\epsilon\mathcal{D}_R\right)
  \mathcal{D}_LV\,,
  \tag{29.5.24}\label{eq:29.5.24}
\end{equation}
where \(V\) is an unrestricted real superfield. Instead, we can implement
the condition \eqref{eq:29.5.23} by introducing a Lagrange multiplier term
in the action
\begin{equation}
  \Delta I_{\mathrm{eff}}
  =\frac{1}{8\pi}\operatorname{Re}\int d^4x\,
    \left[
      \widetilde V
      \left(\mathcal{D}_L^{\mathrm T}\epsilon W_L\right)
    \right]_D,
  \tag{29.5.25}\label{eq:29.5.25}
\end{equation}
where \(\widetilde V\) is an unrestricted real superfield. (The numerical
factor \(1/8\pi\) serves to fix the normalization of \(\widetilde V\) in a
way we shall find convenient later.) Then in path integrals we can integrate
over both \(\widetilde V\) and \(W_L\), with no restrictions on either except
that \(\widetilde V\) is a real superfield and \(W_L\) is a left-handed
spinor superfield satisfying the left-chirality condition
\(\mathcal{D}_{R\alpha}W_{L\beta}=0\). Integrating by parts in superspace
allows us to write the new term in the action as
\begin{equation}
  \Delta I_{\mathrm{eff}}
  =-\frac{1}{8\pi}\operatorname{Re}\int d^4x\,
    \left[
      \left(\mathcal{D}_L\widetilde V\right)^{\mathrm T}
      \epsilon W_L
    \right]_D,
  \tag{29.5.26}\label{eq:29.5.26}
\end{equation}
or, using Eq.~\eqref{eq:26.3.31} and the fact that \(W_L\) is left-chiral,
\begin{equation}
  \Delta I_{\mathrm{eff}}
  =\operatorname{Re}\left[
    \frac{i}{4\pi}\int d^4x\,
    \left(\widetilde W_L^{\mathrm T}\epsilon W_L\right)_{\mathcal F}
  \right],
  \tag{29.5.27}\label{eq:29.5.27}
\end{equation}
where \(\widetilde W_L\) is defined in terms of \(\widetilde V\) by the same
relation that previously gave \(W_L\) in terms of \(V\):
\begin{equation}
  \widetilde W_L
  =\frac{i}{4}
    \left(\mathcal{D}_R^{\mathrm T}\epsilon\mathcal{D}_R\right)
    \mathcal{D}_L\widetilde V.
  \tag{29.5.28}\label{eq:29.5.28}
\end{equation}
But now we are integrating over \(W_L\) with no such constraint, with an
action that is quadratic in \(W_L\)
\begin{equation}
  \begin{aligned}
  I_{\mathrm{eff}}+\Delta I_{\mathrm{eff}}
  ={}&\operatorname{Im}\int d^4x\,
  \left[
    -\frac{1}{8\pi}h'(\Phi)
      \left(W_L^{\mathrm T}\epsilon W_L\right)
    -\frac{1}{4\pi}
      \left(\widetilde W_L^{\mathrm T}\epsilon W_L\right)
  \right]_{\mathcal F}
  \\
  &+\frac{1}{8\pi}\operatorname{Im}\int d^4x\,
    \left[\Phi^*h(\Phi)\right]_D.
  \end{aligned}
  \tag{29.5.29}\label{eq:29.5.29}
\end{equation}
This integration is done by setting \(W_L\) equal to the value where the
action is stationary in \(W_L\):
\begin{equation}
  W_L=-\frac{\widetilde W_L}{h'(\Phi)}\,,
  \tag{29.5.30}\label{eq:29.5.30}
\end{equation}
and the whole effective action becomes
\begin{equation}
  \begin{aligned}
  \widetilde I_{\mathrm{eff}}
  ={}&+\frac{1}{8\pi}\operatorname{Im}\int d^4x\,
    \left[
      \frac{1}{h'(\Phi)}
      \left(\widetilde W_L^{\mathrm T}\epsilon\widetilde W_L\right)
    \right]_{\mathcal F}
  \\
  &+\frac{1}{8\pi}\operatorname{Im}\int d^4x\,
    \left[\Phi^*h(\Phi)\right]_D.
  \end{aligned}
  \tag{29.5.31}\label{eq:29.5.31}
\end{equation}
Now, if we define a new left-chiral scalar superfield and a new \(h\)
function
\begin{equation}
  \widetilde\Phi\equiv h(\Phi)\,,
  \qquad
  \widetilde h(\widetilde\Phi)\equiv-\Phi\,,
  \tag{29.5.32}\label{eq:29.5.32}
\end{equation}
then
\begin{equation}
  \frac{d\widetilde h}{d\widetilde\Phi}\frac{dh}{d\Phi}
  =-\frac{d\Phi}{d\widetilde\Phi}
    \frac{d\widetilde\Phi}{d\Phi}
  =-1\,,
  \tag{29.5.33}\label{eq:29.5.33}
\end{equation}
so the action \eqref{eq:29.5.31} may be written as
\begin{equation}
  \begin{aligned}
  \widetilde I_{\mathrm{eff}}
  ={}&-\frac{1}{8\pi}\operatorname{Im}\int d^4x\,
    \left[
      \widetilde h'(\widetilde\Phi)
      \left(\widetilde W_L^{\mathrm T}\epsilon\widetilde W_L\right)
    \right]_{\mathcal F}
  \\
  &+\frac{1}{8\pi}\operatorname{Im}\int d^4x\,
    \left[\widetilde\Phi^*\widetilde h(\widetilde\Phi)\right]_D.
  \end{aligned}
  \tag{29.5.34}\label{eq:29.5.34}
\end{equation}
From the method by which we have derived it, we know that the theory based on
this effective action is equivalent to the original effective field theory,
so \textit{the \(N=2\) effective field theory with scalar field value \(a\)
and \(h\) function \(h(a)\) is physically equivalent to one with scalar
field value} \(a_D\)\footnote{The subscript \(D\) stands for `dual,' and of
course has nothing to do with the \(D\)-term of a superfield.}
\(\equiv h(a)\) \textit{and \(h\) function}
\(\widetilde h(a_D)\equiv-a\). This is the version of duality that applies
in the present context.

There is another transformation of the function \(h(\Phi)\) that also yields
an equivalent Lagrangian (now with no change in \(\Phi\) or \(W\)) and can be
combined with the transformation \eqref{eq:29.5.32} to give a larger group of
duality transformations. Suppose we shift \(h(\Phi)\) by a linear term with
a real coefficient:
\begin{equation}
  h(\Phi)\longrightarrow h(\Phi)+b\Phi\,,
  \tag{29.5.35}\label{eq:29.5.35}
\end{equation}
where \(b\) is a real constant. Then the first term in the effective
Lagrangian density \eqref{eq:29.5.22} is shifted by
\(b\operatorname{Im}\{[\Phi^*\Phi]_D\}/8\pi\), which vanishes since the
\(D\)-term of \(\Phi^*\Phi\) is real. The change in the effective Lagrangian
density \eqref{eq:29.5.22} is therefore given by the shift in the second term
\begin{equation}
  \mathcal{L}_{\mathrm{eff}}\longrightarrow\mathcal{L}_{\mathrm{eff}}
  -\frac{b}{8\pi}\operatorname{Im}
    \left[\left(W_L^{\mathrm T}\epsilon W_L\right)\right]_{\mathcal F},
  \tag{29.5.36}\label{eq:29.5.36}
\end{equation}
or, according to Eq.~\eqref{eq:27.2.13},
\begin{equation}
  \mathcal{L}_{\mathrm{eff}}\longrightarrow\mathcal{L}_{\mathrm{eff}}
  -\frac{b}{8\pi}\left[
    i\left(\bar\lambda\sl{\partial}\gamma_5\lambda\right)
    +\frac14\epsilon^{\mu\nu\rho\sigma}f_{\mu\nu}f_{\rho\sigma}
  \right].
  \tag{29.5.37}\label{eq:29.5.37}
\end{equation}
The first term in the brackets on the right-hand side is a spacetime
derivative, and therefore does not affect the effective action. The second
term in the brackets can also be written as a spacetime derivative
\((1/2)\partial_\mu(\epsilon^{\mu\nu\rho\sigma}A_\nu f_{\rho\sigma})\)
wherever \(f_{\mu\nu}\) can be written as
\(\partial_\mu A_\nu-\partial_\nu A_\mu\). However, as discussed in Section
23.3, the gauge transformation that we used to put \(\phi_A\) in the
3-direction must be singular somewhere, so we cannot write \(f_{\mu\nu}\)
everywhere in terms of a single \(A_\mu\). In consequence, physical
quantities \textit{can} be affected by a term in the action of the form
\begin{equation}
  -\frac{\theta}{64\pi^2}\int d^4x\,
    \epsilon^{\mu\nu\rho\sigma}f_{\mu\nu}f_{\rho\sigma}.
  \tag{29.5.38}\label{eq:29.5.38}
\end{equation}
In particular, Witten~\hyperref[ch29-ref-11]{[11]} has shown that, in the
theory of magnetic monopoles described in Section 23.3, the electric charge
of a magnetic monopole with minimum monopole moment in the presence of such
a term is \(e\theta/2\pi\). But, as mentioned in Section 27.9, in this theory
there are also dyons, particles with both magnetic monopole moments and
charges given by any integer multiple of \(e\), so the whole pattern of
monopole and dyon charges is periodic in \(\theta\) with period \(2\pi\). In
fact, all physical quantities have this periodicity, because the \(\theta\)
of the low-energy effective theory is inherited from the \(\theta\) appearing
in the Lagrangian density \eqref{eq:29.5.1} of the underlying theory (note
Eq.~\eqref{eq:29.5.21}), and we have seen in Section 23.5 that all physical
quantities are periodic in this \(\theta\). According to
Eq.~\eqref{eq:29.5.37}, the transformation \eqref{eq:29.5.35} changes
\(\theta\) by \(2\pi b\), and therefore with \(b\) equal to an arbitrary
positive or negative integer it yields an equivalent effective action.

This sort of transformation is closely related to an exact invariance of the
underlying theory. The Lagrangian \eqref{eq:29.5.1} is invariant under a
continuous \(R\)-symmetry, under which
\begin{equation}
  \theta_L\longrightarrow\exp(i\alpha)\theta_L\,,
  \qquad
  W_{AL}\longrightarrow\exp(i\alpha)W_{AL}\,,
  \qquad
  \Phi_A\longrightarrow\exp(2i\alpha)\Phi_A.
  \tag{29.5.39}\label{eq:29.5.39}
\end{equation}
This symmetry is broken by an anomaly: both \(\lambda_{AL}\) and
\(\psi_{AL}\) have \(R\) quantum numbers \(+1\), and of course
\(\lambda_{AR}\) and \(\psi_{AR}\) have \(R\) quantum numbers \(-1\), so
according to Eqs.~\eqref{eq:23.5.21} and \eqref{eq:23.5.23}, the measure in
the integral over fermion fields changes under the transformation
\eqref{eq:29.5.39} by a factor \(\exp(2i\alpha N\nu)\), where \(\nu\) is an
integer, the winding number of the vector field configuration, and \(N\) is
defined by
\[
  \operatorname{Tr}(t_At_B)=\frac12N\delta_{AB}.
\]
(The factor \(\exp(2i\alpha N\nu)\) is the same as given in Section 23.5,
even though here we have two fields \(\psi_A\) and \(\lambda_A\) in the same
representation of the gauge group, because these are Majorana fields, in
contrast with the Dirac fields used in Section 23.5.) The generators here
are \((t_A)_{BC}=-i\epsilon_{ABC}\), so \(N=4\), and therefore the measure
remains invariant when \(\exp(8i\alpha)=1\). In other words, the continuous
\(R\)-symmetry is broken by instantons to a \(\mathbb{Z}_8\) subgroup,
generated by the transformation
\begin{equation}
  \psi_{AL}\longrightarrow\sqrt{i}\,\psi_{AL}\,,
  \qquad
  \lambda_{AL}\longrightarrow\sqrt{i}\,\lambda_{AL}\,,
  \qquad
  \phi_A\longrightarrow i\phi_A.
  \tag{29.5.40}\label{eq:29.5.40}
\end{equation}
This symmetry then must carry over to the effective low-energy theory.
However, it does \textit{not} tell us that the effective Lagrangian
\eqref{eq:29.5.20} is invariant under the discrete transformation
\eqref{eq:29.5.40}, which would require that \(-ih(ia)=h(a)\) for all
\(a\). This condition is satisfied by the tree approximation
\eqref{eq:29.5.21}, but (as we will see) it is violated even in one-loop
order. The \(\mathbb{Z}_8\) \(R\)-symmetry is realized in the effective
theory by the condition that the effective theory with \(h\) function
\(-ih(ia)\) is \textit{equivalent} to one with \(h\) function \(h(a)\).
That is, \(-ih(ia)\) must be related to \(h(a)\) by some combination of the
transformations \eqref{eq:29.5.32} and \eqref{eq:29.5.35}, with \(b\) an
integer.

We have seen that the physical significance of the theory is left unchanged
by two sorts of transformation: the transformation
\(\Phi\longrightarrow\widetilde\Phi=h(\Phi)\) and
\(h(\Phi)\longrightarrow\widetilde h(\widetilde\Phi)=-\Phi\), which can be
written as
\begin{equation}
  \begin{pmatrix}
    \Phi\\ h(\Phi)
  \end{pmatrix}
  \longrightarrow
  \begin{pmatrix}
    0&1\\ -1&0
  \end{pmatrix}
  \begin{pmatrix}
    \Phi\\ h(\Phi)
  \end{pmatrix},
  \tag{29.5.41}\label{eq:29.5.41}
\end{equation}
and the transformations
\(h(\Phi)\longrightarrow h(\Phi)+b\Phi\),
\(\Phi\longrightarrow\Phi\), which can be written as
\begin{equation}
  \begin{pmatrix}
    \Phi\\ h(\Phi)
  \end{pmatrix}
  \longrightarrow
  \begin{pmatrix}
    1&0\\ b&1
  \end{pmatrix}
  \begin{pmatrix}
    \Phi\\ h(\Phi)
  \end{pmatrix},
  \tag{29.5.42}\label{eq:29.5.42}
\end{equation}
where \(b\) is an arbitrary integer. By combining these transformations, we
can make generalized duality transformations
\begin{equation}
  \begin{pmatrix}
    \Phi\\ h(\Phi)
  \end{pmatrix}
  \longrightarrow
  \begin{pmatrix}
    n&m\\ k&l
  \end{pmatrix}
  \begin{pmatrix}
    \Phi\\ h(\Phi)
  \end{pmatrix},
  \tag{29.5.43}\label{eq:29.5.43}
\end{equation}
where \(n\), \(m\), \(k\), and \(l\) are any integers satisfying the
condition that the matrix in Eq.~\eqref{eq:29.5.43} should, like those in
Eqs.~\eqref{eq:29.5.41} and \eqref{eq:29.5.42}, have unit determinant:
\begin{equation}
  nl-mk=1.
  \tag{29.5.44}\label{eq:29.5.44}
\end{equation}
These transformations therefore form the group \(SL(2,\mathbb{Z})\).

The physical significance of the duality transformation can be brought
forward by considering the central charge of the \(N=2\) supersymmetry
algebra. As shown in Section 25.5, the eigenvalue \(Z_{12}\) of this central
charge in any one-particle state sets a lower bound
\(M\geq\lvert Z_{12}\rvert/2\) on the mass of the state, a bound that is
reached for particles that belong to `short' supermultiplets. We saw in
Section 27.9 that \(Z_{12}\) is given in the underlying \(SU(2)\) gauge
theory with \(N=2\) supersymmetry by Eq.~\eqref{eq:27.9.22}:
\[
  Z_{12}=2\sqrt{2}\,v[iq-\mathcal{M}]\,,
\]
where \(q\) and \(\mathcal{M}\) are the charge and magnetic moment of the
particle, and \(v\) is the vacuum expectation value of the conventionally
normalized neutral scalar field. In the notation used in the present
section, we have absorbed a factor \(e\) into the normalization of the field
\(a\), so here we have
\begin{equation}
  Z_{12}=2\sqrt{2}\,a[iq-\mathcal{M}]/e.
  \tag{29.5.45}\label{eq:29.5.45}
\end{equation}
This theory contains particles like the massive elementary scalars, spinors,
and vector bosons, which have charge \(e\) and zero magnetic monopole moment,
so such particles are eigenstates of \(Z_{12}\) with eigenvalue
\begin{equation}
  Z_{12}^{\text{charged elementary}}=2\sqrt{2}\,ia.
  \tag{29.5.46}\label{eq:29.5.46}
\end{equation}
As we have seen, the theory with scalar field \(a\) and \(h\) function
\(h(a)\) is equivalent to one with scalar field \(na+mh(a)\) and \(h\)
function \(ka+lh(a)\). Therefore for all integers \(n\) and \(m\) the theory
with scalar field \(a\) and \(h\) function \(h(a)\) must \textit{also}
contain a particle that looks like a massive particle of charge \(e\) and
zero magnetic moment in the version of the theory with scalar field
\(na+mh(a)\) and \(h\) function \(ka+lh(a)\), and that therefore has a
central charge
\begin{equation}
  Z_{12}=2\sqrt{2}\,i[na+mh(a)].
  \tag{29.5.47}\label{eq:29.5.47}
\end{equation}
Comparing this with Eq.~\eqref{eq:29.5.45}, we see that, in the version of
this theory with scalar field \(a\) and \(h\) function \(h(a)\), this
particle appears as one with a charge \(q\) and a magnetic monopole moment
\(\mathcal{M}\) given by
\begin{equation}
  q/e=n+m\operatorname{Re}[h(a)/a]\,,
  \qquad
  \mathcal{M}/e=m\operatorname{Im}[h(a)/a].
  \tag{29.5.48}\label{eq:29.5.48}
\end{equation}
This is a dyon, with both charge and magnetic monopole moment. Note that the
formula \eqref{eq:29.5.48} for the charge of this particle bears out Witten's
earlier result about the charge of magnetic monopoles: adding a term \(ba\)
to the function \(h(a)\) changes the charge of a monopole with \(m=1\) by an
amount \(be=e\Delta\theta/2\pi\).

Using the tree-approximation result \eqref{eq:29.5.21} for \(h(a)\) in
Eq.~\eqref{eq:29.5.48} shows that in this approximation the magnetic
monopole moments in this theory are multiples of a value \(4\pi/e\). This is
the same as the magnetic monopole moment derived in Section 23.3 (recall that
the magnetic monopole moment \(g\) used there is \(\mathcal{M}/4\pi\)), but
this is a semi-classical result, subject to quantum corrections. Note that
the duality transformation \eqref{eq:29.5.41} with \(n=0\) and \(m=1\)
takes an elementary particle with charge \(e\) and magnetic monopole moment
zero into a non-elementary particle with magnetic monopole moment
\(\operatorname{Im}[h(a)/a]\), which in the tree approximation is \(4\pi/e\).

The beta function for the electric charge in the underlying \(N=2\)
supersymmetric \(SU(2)\) gauge theory is given perturbatively by the one-loop
result Eq.~\eqref{eq:27.9.50}, with the first Casimir invariant taken now as
\(C_1=2e\) and the number of hypermultiplets set at \(H=0\):
\begin{equation}
  \beta_{\mathrm{perturbative}}(e)=-\frac{e^3}{4\pi^2}.
  \tag{29.5.49}\label{eq:29.5.49}
\end{equation}
We here take \(a\) as the renormalization scale, so the running charge
\(e(a)\) satisfies
\begin{equation}
  a\frac{d}{da}e(a)=\beta\bigl(e(a)\bigr).
  \tag{29.5.50}\label{eq:29.5.50}
\end{equation}
Using the perturbation theory formula \eqref{eq:29.5.49}, this gives
\begin{equation}
  \left[e^{-2}(a)\right]_{\mathrm{perturbative}}
  =\frac{1}{2\pi^2}\ln\left(\frac{a}{\Lambda}\right),
  \tag{29.5.51}\label{eq:29.5.51}
\end{equation}
where \(\Lambda\) is an integration constant. In the formalism we are
adopting here, with a factor \(e\) absorbed into the definition of the gauge
field, the quantity \(e^{-2}(a)\) appears as the coefficient \(h'(a)/4\pi i\)
of the term \(-\tfrac14f_{\mu\nu}f^{\mu\nu}\) in the low-energy effective
Lagrangian \eqref{eq:29.5.20}, so the function \(h(a)\) is
